\documentclass[12pt]{article}
\usepackage[utf8]{inputenc}
\usepackage[T1]{fontenc}
\usepackage[margin=1in]{geometry}
\usepackage{setspace}
\usepackage{xcolor}
\usepackage[colorlinks=true, urlcolor=blue!60!black, linkcolor=blue!60!black]{hyperref}
\setlength{\parindent}{0em}
\setlength{\parskip}{0.6em}
\renewcommand{\baselinestretch}{1.0}
\begin{document}
\begin{center}
{\LARGE \textit{Barron v.\ Mayor \& City Council of Baltimore}, 32 U.S.\ (7 Pet.) 243 (1833)\par}
\vspace{0.5em}
{\large Case Brief --- Two-Page Edition\par}
\end{center}

\section{Memory Jogger}
A unanimous Court held that the Fifth Amendment's Takings Clause, and by extension the whole Bill of Rights, restrains only the federal government, not the States or their cities (\href{https://supreme.justia.com/cases/federal/us/32/243/}{Justia}).

\section{Detailed Case Facts}
Baltimore's street grading and paving diverted streams toward John Barron's wharf, which had the deepest water in the harbor. Silt washed down by heavy rains made the water too shallow for large vessels, rendering the wharf nearly worthless. The city paid nothing, and Barron claimed a taking without just compensation under the Fifth Amendment.

\section{Procedural History}
A Baltimore County jury awarded Barron \$4,500; Maryland's Court of Appeals for the Western Shore reversed. Barron sought review by writ of error under Section 25 of the Judiciary Act of 1789.

\section{Judicial Votes}
Writ dismissed, 7--0. Chief Justice Marshall wrote for the Court; there were no concurrences or dissents.

\section{Holding}
The Takings Clause, and by the opinion's reasoning the Bill of Rights generally, limits only the federal government. With no federal question presented, the writ was dismissed for want of jurisdiction. Fourteenth Amendment incorporation later largely superseded the rule, starting with takings in 1897.

\section{Analysis of Opinions}
\textbf{Marshall:} The people ordained the Constitution ``for their own government, and not for the government of the individual States.'' Article I, \S~9's general limits bind only Congress, while \S~10 binds the States with express ``No State shall'' language; the Fifth Amendment has no such words. Citizens could amend their own state constitutions rather than use ``[t]he unwieldy and cumbrous machinery'' of Article V. And the amendments answered fears of the new national government, ``not against those of the local governments.''

\section{Examples: Future Applications}
\textbf{Same-Side:} A State may charge murder by information rather than grand jury indictment, since that clause was never incorporated (\textit{Hurtado}). \textbf{Opposite-Side:} A city's \$1 condemnation award violates Fourteenth Amendment due process under \textit{Chicago, Burlington \& Quincy}. \textbf{Fence-Sitter:} State troops quartered in private homes raise the unresolved question whether the Third Amendment is incorporated.

\section{Critique}
Originalists praise \textit{Barron}'s intratextual method, though Crosskey argued only the First Amendment is limited to ``Congress.'' Progressive critics note that Marshall's remedy of state constitutional reform was illusory for the enslaved, free Black citizens, women, and dissenters who lacked political power.

\section{Key Quotations}
\begin{quote}
``The powers they conferred on this government were to be exercised by itself, and the limitations on power, if expressed in general terms, are naturally, and we think necessarily, applicable to the government created by the instrument. They are limitations of power granted in the instrument itself, not of distinct governments framed by different persons and for different purposes.''
\end{quote}
\noindent --- Chief Justice Marshall, Opinion of the Court, 32 U.S. at 247.

\end{document}
